\subsection{分配律公式}

命题8. 设 \(((\mathbf{X}_{\lambda,1})_{i\in \mathbf{J}_2})_{\lambda \in \mathbf{L}}\) 为一个（以L为指标集的）集合族的族。假设 \(\mathbf{L}\neq \emptyset\) 以及 \(J_{\lambda}\neq \emptyset\) 对于每一个 \(\lambda \in \mathbf{L}\) 成立. 设

\[
\mathrm{I} = \prod_ {\lambda \in \mathrm{L}} \mathrm{J} _ {\lambda} \neq \emptyset .
\]

于是我们有

\[
\begin{array}{l} \bigcup_ {\lambda \in \mathbf {L}} \Big (\bigcap_ {t \in \mathbf {J} _ {\lambda}} \mathbf {X} _ {\lambda , t} \Big) = \bigcap_ {f \in \mathbf {I}} \Big (\bigcup_ {\lambda \in \mathbf {L}} \mathbf {X} _ {\lambda , f (\lambda)} \Big), \\ \bigcap_ {\lambda \in \mathbf {L}} \Big (\bigcap_ {t \in \mathbf {J} _ {\lambda}} \mathbf {X} _ {\lambda , t} \Big) = \bigcup_ {f \in \mathbf {I}} \Big (\bigcap_ {\lambda \in \mathbf {L}} \mathbf {X} _ {\lambda , f (\lambda)} \Big) \end{array}
\]

（并集对交集的“分配律”，以及交集对并集的分配律）。

设 \(x\) 为 \(\bigcup_{\lambda \in L} \left( \bigcap_{i \in J_{\lambda}} X_{\lambda, i} \right)\) 的一个元素，并令 \(f\) 为 1 的任意元素。存在一个指标 \(\lambda\) 使得 \(x \in \bigcap_{i \in J_{\lambda}} X_{\lambda, i}\) ；因此 \(x \in X_{\lambda, f(\lambda)}\) ，从而

\[
x \in \bigcup_ {\lambda \in \mathbf {L}} \mathrm{X} _ {\lambda , f (\lambda)}.
\]

由于这对每个 \(f \in I\) ，我们有

\[
x \in \bigcap_ {f \in I} \left(\bigcup_ {\lambda \in L} X _ {\lambda , f (\lambda)}\right).
\]

反之，设 \(x\) 是一个不属于该集合的对象

\[
\bigcup_ {\lambda \in L} \left(\bigcap_ {i \in J _ {\lambda}} X _ {\lambda , i}\right).
\]

于是对每个 \(\lambda \in L\) 我们拥有 \(x \notin \bigcap_{\iota \in J_{\lambda}} X_{\lambda, \iota}\) ，这意味着集合 \(J_{\lambda}'\) 指标 \(\iota \in J_{\lambda}\) 使得 \(x \notin X_{\lambda, \iota}\) 对每个 \(\lambda \in L\) 都不是空的。根据命题5、推论2，存在一个函数图 \(f\) ，其定义域为 \(L\) 使得对于每一个 \(\lambda \in L\) 我们拥有 \(f(\lambda) \in J_{\lambda}'\) 因此 \(f \in I\) 并且 \(x \notin X_{\lambda, f(\lambda)}\) 对于每一个 \(\lambda \in L\) ；因此

\[
x \notin \bigcup_ {\lambda \in \mathbf {L}} \mathrm{X} _ {\lambda , f (\lambda)}
\]

并且从而

\[
x \notin \bigcap_ {f \in I} \left(\bigcup_ {\lambda \in L} X _ {\lambda , f (\lambda)}\right).
\]

这就完成了第一个公式的证明。第二个公式通过将第一个公式应用于族 \(\left((\mathbb{C}_{\Lambda}X_{\lambda,t})_{t\in J_{\lambda}}\right)_{\lambda \in L}\) 而得出，其中 A 表示并集 \(\bigcap_{\lambda \in L}\left(\bigcap_{t\in J_{\lambda}}X_{\lambda}\right)\) .

推论。设 \((\mathbf{X}_i)_{t\in \mathbf{I}}\) 并且 \((\mathbf{Y}_x)_{x\in \mathbb{K}}\) 是两族集合，其指标集 I, K 非空。则

\[
\begin{array}{l} \left(\bigcap_ {i \in I} X _ {i}\right) \cup \left(\bigcap_ {x \in K} Y _ {x}\right) = \bigcap_ {(i, x) \in I \times K} (X _ {i} \cup Y _ {x}), \\ \left(\bigcup_ {i \in I} X _ {i}\right) \cap \left(\bigcup_ {x \in K} Y _ {x}\right) = \bigcup_ {(i, x) \in I \times K} (X _ {i} \cap Y _ {x}). \end{array}
\]

让 \(\alpha, \beta\) 是两个不同的对象；将命题 8 的公式（以 \(L = \{\alpha, \beta\}\) , \(J_{\alpha} = I\) , \(J_{\beta} = K\) ) 到该族 \(((Z_{\lambda, \mu})_{\mu \in J_{\lambda}})_{\lambda \in L}\) ，其中 \(Z_{\alpha, \iota} = X_{\iota}\) 对于所有 \(\iota \in I\) 并且 \(Z_{\beta, x} = Y_{x}\) 对于每一个 \(x \in K\) . 由于典范双射的存在， \(\prod_{\lambda \in L} J_{\lambda}\) 到上 \(I \times K\) (no. 3) 且由§4的命题1，我们得到所述公式。

命题9. 设 \(((\mathbf{X}_{\lambda, t})_{t \in \mathbf{J}_{\lambda}})_{\lambda \in \mathbf{L}}\) 为一族（以指标集 L 为指标）的集合族。设 \(\mathbf{I} = \prod_{\lambda \in \mathbf{J}_{\lambda}} \mathbf{J}_{\lambda}\) 。然后

\[
\prod_ {\lambda \in L} \left(\bigcup_ {i \in J _ {\lambda}} X _ {\lambda , i}\right) = \bigcup_ {f \in I} \left(\prod_ {\lambda \in L} X _ {\lambda , f (\lambda)}\right)
\]

并且（如果 \(\mathbf{L} \neq \phi\) 并且 \(\mathbf{J}_{\lambda} \neq \phi\) 对于每一个 \(\lambda \in \mathbf{L}\) )

\[
\prod_ {\lambda \in \mathrm{L}} \left(\bigcap_ {i \in \mathrm{J} _ {\lambda}} \mathrm{X} _ {\lambda , i}\right) = \bigcap_ {f \in \mathrm{I}} \left(\prod_ {\lambda \in \mathrm{L}} \mathrm{X} _ {\lambda , f (\lambda)}\right)
\]

（乘积对并集和交集的“分配律”）。

第一个公式在以下情况下显然成立：如果 \(L = \emptyset\) 或者如果 \(J_{\lambda} = \emptyset\) 对于某个 \(\lambda \in L\) 。如果不是，设 \(g\) 为 \(\prod_{\lambda \in L} \left( \bigcap_{t \in J_{\lambda}} X_{\lambda, t} \right)\) 对于每一个 \(\lambda \in L\) 存在一个指标 \(t \in J_{\lambda}\) 使得 \(g(\lambda) \in X_{\lambda, t}\) ；换句话说，该集合 \(H_{\lambda}\) 指标 \(t \in J_{\lambda}\) 使得 \(g(\lambda) \in X_{\lambda, t}\) 非空。因此，根据命题5的推论2，存在一个函数图。 \(f\) ，其定义域为 \(L\) 使得

\[
f (\lambda) \in \mathrm{H} _ {\lambda}
\]

对于每一个 \(\lambda \in L\) ，即 \(g(\lambda) \in X_{\lambda, f(\lambda)}\) 。因此我们有 \(g \in \prod_{\lambda \in L} X_{\lambda, f(\lambda)}\) ，从而 \(g \in \bigcup_{f \in I} \left( \prod_{\lambda \in L} X_{\lambda, f(\lambda)} \right)\) 。反之，若

\[
g \in \bigcup_ {f \in \mathbf {I}} \left(\prod_ {\lambda \in \mathbf {L}} \mathrm{X} _ {\lambda , f (\lambda)}\right),
\]

存在一个函数图 \(f \in I\) 使得对于每一个 \(\lambda \in L\) 我们拥有

\[
g (\lambda) \in \mathbf {X} _ {\lambda , f (\lambda)}
\]

并且，更不用说， \(g(\lambda) \in \bigcup_{t \in J_1} X_{\lambda,t}\) 。这就完成了第一个公式的证明。第二个公式的证明类似但更简单，我们留给读者。

推论1。假设 \(\mathbf{L} \neq \emptyset\) 以及 \(J_{\lambda} \neq \emptyset\) 对于每一个 \(\lambda \in L\) . 如果对于 \(e\) achindex \(\lambda \in L\) 族 \((X_{\lambda, t})_{t \in J_{\lambda}}\) 是 \(X_{\lambda} = \bigcup_{t \in J_{\lambda}} X_{\lambda, t}\) 的一个划分，则族 \(\left( \prod_{\lambda \in L} X_{\lambda, f(\lambda)} \right)_{f \in I}\) 是 \(\prod_{\lambda \in L} X_{\lambda}\) 的一个划分.

如果我们设

\[
\mathrm{P} _ {f} = \prod_ {\lambda \in \mathbf {L}} \mathrm{X} _ {\lambda , f (\lambda)},
\]

那么，根据命题9的第一个公式，只需证明 \(\mathbf{P}_f \neq \emptyset\) 对于所有 \(f \in \mathbf{I}\) ，以及 \(\mathbf{P}_f \cap \mathbf{P}_g = \emptyset\) 每当 \(f\) 和 \(g\) 是 \(\mathbf{I}\) 的不同元素。第一点由命题5的推论2得出。至于第二点，如果 \(f \neq g\) , 存在 \(\lambda \in \mathbf{L}\) 使得

\[
f (\lambda) \neq g (\lambda)
\]

因此，根据假设， \(\mathbf{X}_{\lambda,f(\lambda)} \cap \mathbf{X}_{\lambda,g(\lambda)} = \varnothing\) 。由此可知，不存在属于 \(P_{f} \cap P_{g}\) 的图；因为如果G是这样的图，我们就会有 \(\mathbf{G}(\lambda) \in \mathbf{X}_{\lambda,f(\lambda)} \cap \mathbf{X}_{\lambda,g(\lambda)} = \varnothing\) ，这是荒谬的。

推论2。设 \((\mathbf{X}_i)_{i\in \mathbf{I}}\) 并且 \((\mathbf{Y}_x)_{x\in \mathbb{K}}\) 为两个集合族。则

\[
\left(\bigcup_ {i \in I} X _ {i}\right) \times \left(\bigcup_ {x \in K} Y _ {x}\right) = \bigcup_ {(i, x) \in I \times K} (X _ {i} \times Y _ {x})
\]

并且，如果I和K非空，

\[
\left(\bigcap_ {i \in I} X _ {i}\right) \times \left(\bigcap_ {x \in K} Y _ {x}\right) = \bigcap_ {(i, x) \in I \times K} (X _ {i} \times Y _ {x}).
\]

证明遵循命题8推论证明的模式。

命题10。设 \((\mathbf{X}_{t,x})_{(t,x)\in \mathbf{I}\times \mathbf{K}}\) 是一个集合族，其指标集是两个集合 I 和 K 的乘积。如果 \(\mathbf{K}\neq \emptyset\) ，我们有

\[
\bigcap_ {x \in \mathbf {K}} \left(\prod_ {i \in I} X _ {i, x}\right) = \prod_ {i \in I} \left(\bigcap_ {\mathbf {K} \in x} X _ {i, x}\right).
\]

待证等式两边都是函数图。一个图 \(f\) 属于左边当且仅当，对每个 \(x \in \mathbf{K}\) , \(f \in \prod_{i \in I} X_{i,x}\) ；也就是说，当且仅当 \(f(t) \in X_{t,x}\) 对于所有 \((t,x) \in I \times K\) 成立。因为 \(f\) 要属于右侧，必要且充分的条件是 \(f(t) \in \bigcap_{x \in K} X_{t,x}\) 对于每一个 \(t \in I\) ，即 \(f(t) \in X_{t,x}\) 对于每一对 \((t, x) \in I \times K\) 都成立。这就完成了证明。

推论。设 \((\mathbf{X}_t)_{t\in \mathbf{I}}\) 并且 \((\mathbf{Y}_t)_{t\in \mathbf{I}}\) 为两个具有相同指标集的集合族 \(\mathbf{I}\neq \emptyset\) 。然后

\[
\begin{array}{l} \left(\prod_ {i \in I} X _ {i}\right) \cap \left(\prod_ {i \in I} Y _ {i}\right) = \prod_ {i \in I} (X _ {i} \cap Y _ {i}), \\ \left(\bigcap_ {i \in I} X _ {i}\right) \times \left(\bigcap_ {i \in I} Y _ {i}\right) = \bigcap_ {i \in I} (X _ {i} \times Y _ {i}). \end{array}
\]

将命题10应用于K（相应地，I）是由两个不同元素构成的集合的情形。

